% Generated by task tex-libraries from dossiers/lib-arklib.md, sections D.0 to D.6 and E, with
% the summary's reading of the master theorem. Snippets are copied from the dossier.
\begingroup\sloppy
\section{ArkLib: what round-by-round knowledge soundness guarantees}\label{sec:gen-lib-ark-sec}

\begin{define}{What the soundness master theorem means}
\lean{piop_rbrKnowledgeSoundness} says: \emph{for every public input and every transcript, if the
verifier accepts then the stack sent as the first message satisfies \lean{M3Holds}, unless some
challenge of the transcript lies in the bad set of its prefix, and every bad set has probability
at most \lean{piopError} of its round.} The probability is over one uniform challenge at a time,
at every prefix (the worst-case form); no prover appears. The step from this to \enquote{an
interactive prover convinces the verifier without a good first message with probability at most
$\sum$ \lean{piopError}} is a union bound that ArkLib admits at both pins
(\lean{rbrKnowledgeSoundness_}\allowbreak\lean{implies_}\allowbreak\lean{knowledgeSoundness}); the step to Fiat--Shamir security is
stated nowhere in ArkLib. At the oracle-protocol level this is soundness with respect to the
committed oracle; knowledge (extracting the stack from a commitment) belongs to the compiled
protocol, for which ArkLib offers no theorem (\lean{BCSTransform} commented out,
\lean{Commitment.extractability} a placeholder whose body is \lean{False}).
\end{define}

This section derives that reading. Revisions: ArkLib at the old pin \code{dca90385} for every
snippet, leanerVM at \code{b435631}; every probe ran at the old pin and none could be re-run after
the upgrade (\cref{ch:probes}).

\subsection{The statement under review}\label{sec:gen-lib-ark-statement}

leanerVM's soundness master theorem is an instance of one ArkLib predicate,
\lean{Verifier.}\allowbreak\lean{rbrKnowledgeSoundnessWorstCaseWith}, applied to the verifier of the composed oracle
protocol after it has been turned into a plain verifier (\lean{toVerifier},
\cref{sec:gen-lib-ark-toverifier}):
\begin{leancode}
/-- **Round-by-round knowledge soundness** at `piopError P`, for the extractor
`piopExtractor P S` (the commit phase's, then the phases') and its knowledge state function,
composed from the phases' own by a proved composition: each challenge can turn the state from
false to true with probability at most its error. The extractor's stack satisfies `M3Holds` (see
the module docstring for the plain reading and what it waits on). -/
theorem piop_rbrKnowledgeSoundness (P : Phases I) (S : P.Security) {σ : Type}
    (init : ProbComp σ) (impl : QueryImpl []ₒ (StateT σ ProbComp)) :
    (leanVmVerifier P).toVerifier.rbrKnowledgeSoundnessWorstCaseWith init impl (M3Rel I)
      (Seam.done I) S.toDef.witMid (piopExtractor P S) (S.toDef.kSF init impl)
      (piopError P) :=
  S.toDef.rbr init impl
\end{leancode}
\src{\file{LeanerVM/}\allowbreak\file{Protocol/}\allowbreak\file{Spine/}\allowbreak\file{Compose.lean:167-177}, leanerVM \code{b435631}}

Its arguments are: the input relation \lean{M3Rel I} (pairs of a public input and a stack
\lean{q}), the output relation \lean{Seam.done I} (everything), the family of intermediate
witness types \lean{S.toDef.witMid}, the extractor \lean{piopExtractor P S}, the knowledge state
function \lean{S.toDef.kSF init impl}, and the error per challenge \lean{piopError P}. The shared
oracle is empty (\lean{[]ₒ}).

\subsection{The definition, in mathematical language}\label{sec:gen-lib-ark-definition}

\paragraph{Setting.} A protocol has $n$ rounds. Each round is either a prover message of a type
$M_k$ or a verifier challenge of a finite type $C_k$. A \emph{prefix} $\tau \in \mathrm{Tr}_k$ is
the list of the first $k$ entries (messages and challenges); $\mathrm{Tr}_n$ is the set of full
transcripts. Statements are $s \in S$, input witnesses $W_{\mathrm{in}}$, output statements $T$,
output witnesses $W_{\mathrm{out}}$. The relations are
$R_{\mathrm{in}} \subseteq S \times W_{\mathrm{in}}$ and
$R_{\mathrm{out}} \subseteq T \times W_{\mathrm{out}}$. The verifier is a function
$V : S \times \mathrm{Tr}_n \to$ (computation with failure) $T$; \emph{accepting} means returning
some $t \in T$, and \emph{rejecting} means failing.

\begin{define}{The round-by-round extractor (\lean{Extractor.RoundByRound})}
A family of intermediate witness types $W_0, \dots, W_n$ with $W_0 = W_{\mathrm{in}}$, and
functions
\[
  \mathit{extractOut} : S \times \mathrm{Tr}_n \times W_{\mathrm{out}} \to W_n,
  \qquad
  \mathit{extractMid}_k : S \times \mathrm{Tr}_{k+1} \times W_{k+1} \to W_k
  \quad (k < n),
\]
from the output witness to the last intermediate witness, and one round backwards. It has no law
of its own. It is given no query log, no access to the shared oracle and no randomness, and
nothing bounds its running time.
\end{define}
\begin{leancode}
structure RoundByRound
    (oSpec : OracleSpec ι) (StmtIn WitIn WitOut : Type) {n : ℕ} (pSpec : ProtocolSpec n)
    (WitMid : Fin (n + 1) → Type) where
  /-- The first intermediate witness type is equal to the input witness type -/
  eqIn : WitMid 0 = WitIn
  /-- Extract intermediate witness for round `m` from intermediate witness for round `m+1`,
    using the transcript up to round `m+1` -/
  extractMid : (m : Fin n) → StmtIn → Transcript m.succ pSpec → WitMid m.succ → WitMid m.castSucc
  /-- Construct the intermediate witness for the final round from the output witness -/
  extractOut : StmtIn → FullTranscript pSpec → WitOut → WitMid (.last n)
\end{leancode}
\src{\file{ArkLib/}\allowbreak\file{OracleReduction/}\allowbreak\file{Security/}\allowbreak\file{RoundByRound.lean:77-86}, ArkLib \code{dca90385}}

\begin{define}{The knowledge state function (\lean{Verifier.}\allowbreak\lean{KnowledgeStateFunction})}
A family of predicates $K_k \subseteq S \times \mathrm{Tr}_k \times W_k$, $k = 0, \dots, n$, with
three laws.
\begin{enumerate}
\item \emph{Start} (\lean{toFun_empty}): for all $s$ and $w \in W_0$,
  \[ (s, w) \in R_{\mathrm{in}} \iff K_0(s, \emptyset, w). \]
\item \emph{Prover rounds} (\lean{toFun_next}): for every round $k$ that is a prover message, all
  $s$, $\tau \in \mathrm{Tr}_k$, messages $m \in M_k$ and $w \in W_{k+1}$,
  \[ K_{k+1}(s, \tau \| m, w) \implies K_k\bigl(s, \tau, \mathit{extractMid}_k(s, \tau\|m, w)\bigr). \]
  No prover message can make the state true unless the witness extracted one round back was
  already good.
\item \emph{End} (\lean{toFun_full}): for all $s$, full $\tau$ and $w_{\mathrm{out}}$,
  \[ \Pr\bigl[V(s,\tau) \text{ returns some } t \text{ with } (t, w_{\mathrm{out}}) \in R_{\mathrm{out}}\bigr] > 0
     \implies K_n\bigl(s, \tau, \mathit{extractOut}(s, \tau, w_{\mathrm{out}})\bigr). \]
  The probability is over the initial state of the shared oracle (drawn by \lean{init}) and the
  oracle's answers (\lean{impl}); for a verifier that makes no query to a shared oracle it is $0$
  or $1$.
\end{enumerate}
\end{define}
\begin{leancode}
structure KnowledgeStateFunction
    (relIn : Set (StmtIn × WitIn)) (relOut : Set (StmtOut × WitOut))
    (verifier : Verifier oSpec StmtIn StmtOut pSpec)
    {WitMid : Fin (n + 1) → Type}
    (extractor : Extractor.RoundByRound oSpec StmtIn WitIn WitOut pSpec WitMid)
    where
  /-- The knowledge state function: takes in round index, input statement, transcript up to that
      round, and intermediate witness of that round, and returns True/False. -/
  toFun : (m : Fin (n + 1)) → StmtIn → Transcript m pSpec → WitMid m → Prop
  /-- The input statement and witness are in the input relation if and only if the state function is
      true for the empty transcript and the input witness -/
  toFun_empty : ∀ stmtIn witMid,
    ⟨stmtIn, cast extractor.eqIn witMid⟩ ∈ relIn ↔ toFun 0 stmtIn default witMid
  /-- If the state function is true for a partial transcript extended with a prover message, then
    the state function is also true for the original partial transcript with the extracted
    intermediate witness -/
  toFun_next : ∀ m, pSpec.dir m = .P_to_V →
    ∀ stmtIn tr msg witMid, toFun m.succ stmtIn (tr.concat msg) witMid →
      toFun m.castSucc stmtIn tr (extractor.extractMid m stmtIn (tr.concat msg) witMid)
  /-- If the verifier can output a statement `stmtOut` that is in the output relation with some
    output witness `witOut`, then the state function is true for the full transcript and the
    extracted last middle witness. -/
  toFun_full : ∀ stmtIn tr witOut,
    Pr[fun stmtOut => (stmtOut, witOut) ∈ relOut
    | OptionT.mk do (simulateQ impl (verifier.run stmtIn tr)).run' (← init)] > 0 →
    toFun (.last n) stmtIn tr (extractor.extractOut stmtIn tr witOut)
\end{leancode}
\src{\file{ArkLib/}\allowbreak\file{OracleReduction/}\allowbreak\file{Security/}\allowbreak\file{RoundByRound.lean:164-189}, ArkLib \code{dca90385}}

\begin{define}{The bound (\lean{Verifier.}\allowbreak\lean{rbrKnowledgeSoundnessWorstCaseWith})}
For a statement $s$, a verifier round $k$ and a prefix $\tau \in \mathrm{Tr}_k$, call \emph{bad
set} the set of challenges
\[
  \mathrm{Bad}_k(s, \tau) = \bigl\{\, c \in C_k : \exists\, w \in W_{k+1},\
  K_{k+1}(s, \tau\|c, w) \wedge \neg K_k\bigl(s, \tau, \mathit{extractMid}_k(s, \tau\|c, w)\bigr) \,\bigr\}.
\]
The predicate says: for \textbf{every} statement $s$, \textbf{every} verifier round $k$ and
\textbf{every} prefix $\tau \in \mathrm{Tr}_k$,
\[
  \Pr_{c \leftarrow C_k \text{ uniform}}\bigl[\, c \in \mathrm{Bad}_k(s, \tau) \,\bigr] \le \varepsilon_k .
\]
\end{define}
\begin{leancode}
/-- Worst-case-per-prefix RBR knowledge soundness for one exact
intermediate-witness family, extractor, and knowledge-state function. -/
def rbrKnowledgeSoundnessWorstCaseWith
    (relIn : Set (StmtIn × WitIn)) (relOut : Set (StmtOut × WitOut))
    (verifier : Verifier oSpec StmtIn StmtOut pSpec)
    (WitMid : Fin (n + 1) → Type)
    (extractor : Extractor.RoundByRound oSpec StmtIn WitIn WitOut pSpec WitMid)
    (kSF : verifier.KnowledgeStateFunction init impl relIn relOut extractor)
    (rbrKnowledgeError : pSpec.ChallengeIdx → ℝ≥0) : Prop :=
  ∀ stmtIn : StmtIn,
  ∀ i : pSpec.ChallengeIdx,
  ∀ transcript : Transcript i.1.castSucc pSpec,
    Pr[fun challenge =>
      ∃ witMid,
        ¬ kSF i.1.castSucc stmtIn transcript
          (extractor.extractMid i.1 stmtIn (transcript.concat challenge) witMid) ∧
          kSF i.1.succ stmtIn (transcript.concat challenge) witMid
      | $ᵗ (pSpec.Challenge i)] ≤ rbrKnowledgeError i
\end{leancode}
\src{\file{ArkLib/}\allowbreak\file{OracleReduction/}\allowbreak\file{Security/}\allowbreak\file{RoundByRound.lean:551-568}, ArkLib \code{dca90385}}

\paragraph{What is quantified over, and what is not.}

{\small
\begin{longtable}{@{}p{0.27\linewidth}p{0.69\linewidth}@{}}
\toprule
Object & Quantified how \\
\midrule
\endhead
\bottomrule
\endfoot
Statements $s$ & universally, including the statements that have a valid witness \\
Verifier rounds $k$ & universally; each has its own error $\varepsilon_k$ \\
Prefixes $\tau$ & universally: every list of prover messages and earlier challenges, whether or
  not any prover could reach it with noticeable probability \\
Intermediate witnesses $w$ & existentially inside the event: the worst one counts \\
The challenge $c$ & the only random variable: uniform on $C_k$ (the \lean{SampleableType}
  instance; its two laws, full support and equal probabilities, force the uniform distribution on
  a finite type, \file{VCVio/}\allowbreak\file{OracleComp/}\allowbreak\file{Constructions/}\allowbreak\file{SampleableType.lean:44-47}) \\
Provers & \textbf{not at all}: no prover, no prover state and no input witness appears \\
The shared oracle & only inside law 3 (can the verifier accept). The state function and the
  extractor cannot depend on the oracle's state or on a query log \\
\lean{WitMid}, the extractor, the state function & named arguments of the predicate
  (\enquote{With}); in \lean{rbrKnowledgeSoundnessWorstCase} they are existentially quantified \\
\end{longtable}}

\paragraph{What \enquote{worst case} adds.} ArkLib's older predicate \lean{rbrKnowledgeSoundness}
bounds the same event in a game: a prover is run against the challenger up to round $k$, which
produces the prefix, then the challenge is drawn, and the bound is on the probability over the
whole game (so it is an average over the earlier challenges). It quantifies over provers, and the
prefix is a random variable:
\begin{leancode}
def rbrKnowledgeSoundness (relIn : Set (StmtIn × WitIn)) (relOut : Set (StmtOut × WitOut))
    (verifier : Verifier oSpec StmtIn StmtOut pSpec)
    (rbrKnowledgeError : pSpec.ChallengeIdx → ℝ≥0) : Prop :=
  ∃ WitMid : Fin (n + 1) → Type,
  ∃ extractor : Extractor.RoundByRound oSpec StmtIn WitIn WitOut pSpec WitMid,
  ∃ kSF : verifier.KnowledgeStateFunction init impl relIn relOut extractor,
  ∀ stmtIn : StmtIn,
  ∀ witIn : WitIn,
  ∀ prover : Prover oSpec StmtIn WitIn StmtOut WitOut pSpec,
  ∀ i : pSpec.ChallengeIdx,
    Pr[fun ⟨transcript, challenge, _proveQueryLog⟩ =>
      ∃ witMid,
        ¬ kSF i.1.castSucc stmtIn transcript
          (extractor.extractMid i.1 stmtIn (transcript.concat challenge) witMid) ∧
          kSF i.1.succ stmtIn (transcript.concat challenge) witMid
    | do
      (simulateQ (impl.addLift challengeQueryImpl : QueryImpl _ (StateT σ ProbComp))
        (do
          let ⟨⟨transcript, _⟩, proveQueryLog⟩ ← prover.runWithLogToRound i.1.castSucc stmtIn witIn
          let challenge ← liftComp (pSpec.getChallenge i) _
          return (transcript, challenge, proveQueryLog))).run' (← init)] ≤
      rbrKnowledgeError i
\end{leancode}
\src{\file{ArkLib/}\allowbreak\file{OracleReduction/}\allowbreak\file{Security/}\allowbreak\file{RoundByRound.lean:416-437}, ArkLib \code{dca90385}}

The worst-case form implies the averaged one with the same errors (proved:
\lean{Verifier.}\allowbreak\lean{rbrKnowledgeSoundnessWorstCaseWith_}\allowbreak\lean{implies_}\allowbreak\lean{rbrKnowledgeSoundnessWith},
\file{Security/}\allowbreak\file{RoundByRound.lean:625}, three standard axioms, probe \code{AxiomsArkLibBuilt}). The
converse fails: from the averaged bound at round $k$ one gets, for one fixed value of the earlier
challenges, only $\varepsilon_k$ multiplied by the inverse probability of those challenges. The
difference matters exactly where it should: a Fiat--Shamir (state-restoration) prover chooses
among many prefixes, so it needs the bound at every prefix, which only the worst-case form gives.
The worst-case form is the one of the literature (\cref{sec:gen-lib-ark-literature}), and the one
leanerVM states.

\paragraph{The plain reading, proved by probe.} Run the extractor backwards along a full
transcript $\tau$:
\[
  w_n = \mathit{extractOut}(s, \tau, w_{\mathrm{out}}), \qquad
  w_k = \mathit{extractMid}_k(s, \tau_{\le k}, w_{k+1}), \qquad
  \mathrm{Ext}(s, \tau, w_{\mathrm{out}}) = w_0 \in W_{\mathrm{in}}.
\]
Then, for any verifier, extractor and knowledge state function, with \textbf{no} hypothesis on the
errors: for all $s$, $\tau$, $w_{\mathrm{out}}$,
\[
\begin{aligned}
  &\Pr\bigl[V(s,\tau) \text{ returns } t,\ (t, w_{\mathrm{out}}) \in R_{\mathrm{out}}\bigr] > 0 \\
  &\qquad\implies
  \bigl(s, \mathrm{Ext}(s, \tau, w_{\mathrm{out}})\bigr) \in R_{\mathrm{in}}
  \ \vee\ \exists\, k \text{ a verifier round},\ \tau_k \in \mathrm{Bad}_k(s, \tau_{<k}).
\end{aligned}
\]
This is \lean{Probe.}\allowbreak\lean{accept_}\allowbreak\lean{imp_}\allowbreak\lean{extract_}\allowbreak\lean{or_bad} of the probe \code{PlainReading}
(\cref{ch:probes}), proved in 40 lines by induction on the round, on the three standard axioms.
With the bound of the predicate, each bad set has probability at most $\varepsilon_k$. So the
predicate reduces knowledge soundness to: \enquote{no challenge of the accepted transcript fell in
the bad set of its prefix}. What remains to get a statement about a prover is a union bound over
the verifier rounds (interactive prover: $\sum_k \varepsilon_k$) or over the prover's queries
(Fiat--Shamir prover with $Q$ queries: about $Q \cdot \max_k \varepsilon_k$). That last step is
what ArkLib admits (\cref{sec:gen-lib-ark-implications}); it is a statement about ArkLib's
execution semantics, not about the protocol.

\paragraph{Instantiated to leanVM.} $S$ is the public input (there is no input oracle),
$W_{\mathrm{in}}$ is the stack $q$, the first message is the stack, $W_{\mathrm{out}}$ is trivial
and $R_{\mathrm{out}}$ is everything. The commit phase's extractor returns the first message
whatever intermediate witness it is handed (\lean{Component.sendExtractor},
\file{LeanerVM/}\allowbreak\file{Protocol/}\allowbreak\file{ToArkLib/}\allowbreak\file{SendOracle.lean:134-138}), so
$\mathrm{Ext}(s, \tau, ()) = \tau_0$. The master theorem therefore reads as in the box at the head
of this section. At the level of the oracle protocol this is a soundness statement about the
committed oracle. The extraction of the stack from a commitment is not in it: it belongs to the
compiled protocol (WHIR, Merkle trees, the random oracle), which the blueprint places in Layers 11
and 12 behind assumed interfaces.

\subsection{Comparison with the literature}\label{sec:gen-lib-ark-literature}

\paragraph{Round-by-round soundness.} Block, Garreta, Katz, Thaler, Tiwari and Zając
(\emph{Fiat--Shamir Security of FRI and Related SNARKs}, ePrint 2023/1071, Definition 3.12, read
from the paper's PDF; it says it follows [CMS19], which follows Canetti et al.\ 2019) define it
with a \enquote{doomed set} $D$ of partial and complete transcripts (a state function taking the
value \emph{doomed}):
\begin{quote}\small
1. If $x \notin L$ then $(x, \emptyset) \in D$. 2. For any complete transcript $\tau$, if
$(x, \tau) \in D$ then $V(x, \tau) = \mathrm{reject}$. 3. If $(x, \tau)$ is an $(i-1)$-round
partial transcript with $(x, \tau) \in D$, then for every potential prover next message $m$,
$\Pr_{c \leftarrow C_i}[(x, \tau, m, c) \notin D] \le \varepsilon(i)$.
\end{quote}
ArkLib's \lean{Verifier.StateFunction} (\file{Security/}\allowbreak\file{RoundByRound.lean:135-152}) is this notion
for a \emph{reduction} (the final condition is \enquote{the verifier does not output a statement
of the output language}) with two differences of shape. Condition 3 is split into a law for prover
messages (\lean{toFun_next}: a doomed prefix stays doomed after any prover message) and the
challenge bound, which lets ArkLib schedule messages and challenges in any order rather than in
alternation; the two forms are equivalent for alternating protocols (extend the doomed set to a
prefix ending in a prover message by the value at the prefix before it). Condition 1 is an
equivalence, \lean{stmt ∈ langIn ↔ toFun 0 stmt default}; the extra direction is harmless (a state
function may always be defined so) and is what sequential composition uses. Canetti et al.'s
original definition (as the dossier's author recalls it; not re-read) is the same three conditions
with a function $\mathrm{State}$ in place of the set $D$.

\paragraph{Round-by-round knowledge soundness.} The same paper, Definition 3.13 (verbatim):
\begin{quote}\small
A (public-coin) holographic IOP $\Pi$ for an indexed relation $R$ has round-by-round knowledge
with error $\varepsilon_k$ if there exists a polynomial-time extractor Ext and for all $i$ there
exists a (not necessarily efficiently computable) \enquote{doomed set} $D(i)$ of partial and
complete transcripts such that: \textbullet\ $(x, \emptyset) \in D(i)$ for all possible input
$x$, regardless of whether $x \in L$ or not. \textbullet\ For any possible input $x$ and any
complete transcript $(x, \tau)$, if $(x, \tau) \in D(i)$ then $V(x, \tau) = \mathrm{reject}$.
\textbullet\ If $i \in [\mu]$ and $(x, \tau)$ is an $(i-1)$-round partial transcript such that
$(x, \tau) \in D(i)$, then for every potential prover next message $m$, it holds that, if
$\Pr_{c \leftarrow C_i}[(x, \tau, m, c) \notin D(i)] > \varepsilon_k(i)$, then
$\mathrm{Ext}(i, x, \tau, m)$ outputs a valid witness for $x$.
\end{quote}
Chiesa and Yogev's book gives (as the dossier's author recalls it) the same notion with a state
function: whenever a prefix is doomed and the next challenge leaves the doomed set with
probability above the error, the extractor, given the prefix up to and including the prover's last
message, outputs a witness. ArkLib's own docstring says that its knowledge state function
\enquote{deliberately differs from ABF26 Definition A.5} (Arnon, Boneh, Fenzi, \emph{Open Problems
in List Decoding and Correlated Agreement}, 2026, per ArkLib's bibliography; not read for the
dossier), so ArkLib models a witness-carrying state function of that manuscript, not
Definition 3.13 directly.

\paragraph{Where ArkLib's definition differs, and whether it weakens the guarantee.}
\begin{enumerate}
\item \emph{The extractor's resources.} Definition 3.13 asks for a polynomial-time extractor.
  ArkLib's \lean{Extractor.RoundByRound} is any Lean function: classical choice is allowed
  (\cref{sec:gen-lib-ark-extractors}, probe \code{Extractors}), and nothing bounds running time.
  In the existential forms (\lean{rbrKnowledgeSoundness}, \lean{rbrKnowledgeSoundnessWorstCase},
  and leanerVM's \lean{piop_}\allowbreak\lean{rbrKnowledgeSoundness_}\allowbreak\lean{exists}) this removes the knowledge content: the
  probe proves \enquote{knowledge soundness at error 0} for an arbitrary relation from a verifier
  that merely decides the language, with an extractor that chooses a witness. \textbf{This is a
  real weakening of the notion as a definition.} It is repaired only by the named form together
  with a reading of the named extractor (what leanerVM does: \lean{piopExtractor} returns the
  first message). ArkLib's own conversion \lean{Extractor.}\allowbreak\lean{RoundByRoundOneShot.}\allowbreak\lean{toRoundByRoundOfRel}
  (\file{Security/}\allowbreak\file{RoundByRound.lean:102-123}) uses \lean{h.choose}, and its docstring says why:
  \enquote{\lean{extractMid} is a mathematical function rather than an algorithm with a tracked
  running time}.
\item \emph{What the extractor sees and when.} Definition 3.13 extracts from the prefix
  $(\tau, m)$ at the round where the prover becomes likely to leave the doomed set. ArkLib
  extracts backwards from the full transcript and the output witness. For an IOP whose witness is
  its first oracle message (leanVM), both read that message, and ArkLib's chain is the identity on
  it. For the Fiat--Shamir reading, ArkLib's form is at least as usable: the extractor is run
  once, on the final accepting transcript; the bound on each prefix's bad set gives the union
  bound over the prover's queries. This is the dossier's analysis; ArkLib proves no Fiat--Shamir
  theorem (\cref{sec:gen-lib-ark-implications}).
\item \emph{The state is on (transcript, intermediate witness).} Definition 3.13's doomed set is on
  transcripts alone and the empty transcript is always doomed; ArkLib's state at the empty
  transcript is \enquote{the witness is valid} and the bad event has an existential over
  intermediate witnesses. When all intermediate witnesses are trivial (every leanVM phase after the
  commit has \lean{witMid := fun _ ↦ Unit}), ArkLib's knowledge state function is a transcript
  predicate and the notion is round-by-round \emph{soundness} of the verifier with respect to the
  seams, which is exactly what a phase over an already committed oracle can offer.
\item \emph{Errors.} ArkLib's error \lean{pSpec.ChallengeIdx → ℝ≥0} may depend on the protocol (in
  leanVM on the instance \lean{I}, which carries the sizes) but not on the statement;
  Definition 3.13's $\varepsilon_k(i)$ depends on the index. Equivalent for leanVM, whose sizes
  index the protocol family.
\item \emph{Acceptance.} ArkLib's law 3 triggers on any positive probability of a related output,
  so a doomed state must reject with probability 1, as in Definition 3.13's second item.
\item \emph{No prover appears} in the worst-case form; Definition 3.13 has none either (\enquote{for
  every potential prover next message}). The averaged \lean{rbrKnowledgeSoundness} does quantify
  over provers and is the weaker notion.
\end{enumerate}

\paragraph{State-restoration knowledge soundness and straight-line extraction.} State-restoration
soundness (Ben-Sasson, Chiesa, Spooner 2016; Chiesa--Yogev) lets the prover rewind the verifier to
any earlier state at most $Q$ times; round-by-round soundness implies it with error about
$Q \cdot \varepsilon$ (Theorem 3.15 of the paper above, quoted in
\cref{sec:gen-lib-ark-implications}). ArkLib defines the state-restoration games
(\file{Security/}\allowbreak\file{StateRestoration.lean:133-154} at the pin) and states no theorem from
round-by-round to state restoration (the statements are commented out,
\file{Security/}\allowbreak\file{Implications.lean:230-254}); the theorems from state restoration to plain soundness
have \lean{sorry} inside their statements (\file{:269}, \file{:284}). Straight-line extraction in
the random oracle model (the extractor reads the prover's oracle queries and never rewinds) is the
shape of ArkLib's \lean{Extractor.Straightline} (\file{Security/}\allowbreak\file{Basic.lean:248-254}, which takes
the prover's and the verifier's query logs), but no ArkLib theorem produces such an extractor from
a round-by-round one at the pin, and the round-by-round extractor itself receives no query log.

\subsection{Which standard consequences are proved and which are admitted}\label{sec:gen-lib-ark-implications}

Status by \texttt{\#print axioms} (probe \code{AxiomsArkLibBuilt}) for the modules in leanerVM's
build; by the source and ArkLib's committed axiom baseline (\cref{sec:gen-lib-ark-baseline}) for
the others (\file{Security/}\allowbreak\file{Implications.lean} is not in leanerVM's build). Lines are those of
\code{dca90385}.

{\footnotesize
\begin{longtable}{@{}p{0.44\linewidth}p{0.20\linewidth}p{0.17\linewidth}p{0.11\linewidth}@{}}
\caption{The standard implications between ArkLib's notions, at the two pins.}\label{tab:gen-lib-ark-implications}\\
\toprule
Implication & Where & At \code{dca90385} & At \code{7653a901} \\
\midrule
\endfirsthead
\toprule
Implication & Where & At \code{dca90385} & At \code{7653a901} \\
\midrule
\endhead
\bottomrule
\endfoot
worst-case rbr knowledge $\Rightarrow$ averaged rbr knowledge
  (\lean{rbrKnowledgeSoundnessWorstCase_}\allowbreak\lean{implies_}\allowbreak\lean{rbrKnowledgeSoundness},
  \lean{…With_implies_…With})
  & \file{Security/}\allowbreak\file{RoundByRound.lean:606, 625} & proved & proved \\
worst-case rbr soundness $\Rightarrow$ averaged rbr soundness & \file{:589} & proved & proved \\
rbr knowledge $\Rightarrow$ rbr soundness
  (\lean{rbrKnowledgeSoundness_}\allowbreak\lean{implies_rbrSoundness})
  & \file{Security/}\allowbreak\file{Implications.lean:87} & proved (not in the baseline; the
    proof is in the file) & proved \\
one-shot rbr knowledge $\Rightarrow$ rbr knowledge
  & \file{Security/}\allowbreak\file{RoundByRound.lean:657}
  & proved, with a classical-choice extractor & proved \\
rbr knowledge $\Rightarrow$ plain knowledge soundness
  (\lean{rbrKnowledgeSoundness_}\allowbreak\lean{implies_knowledgeSoundness}, error
  $\sum \varepsilon$)
  & \file{Security/}\allowbreak\file{Implications.lean:223-228} & \textbf{admitted}
    (\lean{by sorry}) & admitted \\
rbr soundness $\Rightarrow$ soundness (\lean{rbrSoundness_}\allowbreak\lean{implies_}\allowbreak\lean{soundness}) & \file{:79-83}
  & \textbf{admitted} & admitted \\
knowledge soundness $\Rightarrow$ soundness & \file{:49-63} & \textbf{admitted} & admitted \\
rbr $\Rightarrow$ state restoration & \file{:230-254} & commented out & commented out \\
state restoration $\Rightarrow$ plain & \file{:259-313} & admitted, \lean{sorry} in the
  statements & same \\
Fiat--Shamir completeness (\lean{fiatShamir_completeness})
  & \file{FiatShamir/}\allowbreak\file{Basic.lean:163-173}
  & \textbf{admitted}, and stated with the challenge oracle fixed to the constant function
    \lean{default} & admitted \\
Fiat--Shamir (knowledge) soundness & --- & not stated (\enquote{TODO}, \file{:175})
  & not stated \\
BCS transform & \file{BCS/Basic.lean} & commented out & commented out \\
\end{longtable}}

So a reader of \lean{piop_rbrKnowledgeSoundness} trusts the \emph{definition}
(\cref{sec:gen-lib-ark-definition}) and, for the plain reading, the union bound. The
transcript-level half of that reading is proved by the probe \code{PlainReading}
(\lean{accept_imp_extract_or_bad}); the probabilistic half (a union bound over the rounds of an
interactive execution, or over the queries of a Fiat--Shamir prover) is the admitted
\lean{rbrKnowledgeSoundness_}\allowbreak\lean{implies_}\allowbreak\lean{knowledgeSoundness} and the absent Fiat--Shamir theorem. For
the expected form of the latter, the paper above states (Theorem 3.15, verbatim): \enquote{if
$(P, V)$ is a public-coin IOP for $R$ with \dots\ round-by-round knowledge error
$\varepsilon_{\mathrm{rbr}\text{-}k}(x)$ \dots\ then $\mathrm{BCS}^H(P, V)$ is a non-interactive
random oracle proof system for $R$ with \dots\ adaptive knowledge error
$\varepsilon_{\mathrm{fs}\text{-}k}(x, Q, \kappa) = Q\cdot\varepsilon_{\mathrm{rbr}\text{-}k}(x)
+ 3(Q^2+1)/2^{\kappa}$ against $Q$-query adversaries}. The blueprint's
\lean{FiatShamirSecurity} interface names this shape (\enquote{\code{error Q · max_i ε_i}},
\file{docs/}\allowbreak\file{roadmap/}\allowbreak\file{protocol-blueprint.md:1224}); nothing in ArkLib at either pin states it.

\subsection{Non-vacuity of the notion}\label{sec:gen-lib-ark-nonvacuity}

The probe \code{NonVacuity} (\cref{ch:probes}) takes the relation \enquote{the statement is a bit
and only \lean{true} has a witness} (\lean{relIn := {p | p.1 = true}}, \lean{relOut := Set.univ}),
for any shared oracle, any \lean{init} and any \lean{impl}.
\begin{itemize}
\item \textbf{No round, accept everything.} \lean{acceptAll0} (\lean{verify := fun _}\allowbreak\lean{ _ => pure ()})
  has \emph{no} knowledge state function at all, for any extractor
  (\lean{acceptAll0_}\allowbreak\lean{no_stateFunction}): law 3 forces the state at \lean{false}, law 1 then puts
  \lean{false} in the language.
\item \textbf{One challenge, accept everything.} \lean{acceptAll1_}\allowbreak\lean{not_}\allowbreak\lean{sound_}\allowbreak\lean{at_zero}: for
  \emph{every} extractor and \emph{every} knowledge state function, the verifier is not worst-case
  round-by-round knowledge sound at error $0$. The attempt fails exactly where it should: at the
  statement \lean{false}, the empty prefix and the challenge \lean{true}, the state after the
  challenge is forced true by acceptance and the state before is forced false by the relation, so
  the bad set is everything and its probability is $1$, not $0$.
\item \textbf{The error carries the content.} \lean{acceptAll1_sound_at_one}: the same verifier
  \emph{is} knowledge sound at error $1$ (state \enquote{valid witness} at round $0$, \lean{True}
  afterwards). Every verifier with a challenge satisfies the predicate at error $1$. So a theorem
  of the form \lean{rbrKnowledgeSoundnessWorstCaseWith … ε} says nothing until $\varepsilon$ is
  read; in leanerVM the error \lean{piopError P := P.toDef.err} is \emph{data supplied by each
  phase's definition} (\lean{Component.Def.err}), not derived from anything, and a phase may charge
  $1$ to a challenge and discharge its \lean{rbr} obligation trivially. The blueprint's planned
  \lean{piopError_le} (\lean{Σ i, piopError s i ≤ 2^40/|E| + flockError},
  \file{docs/}\allowbreak\file{roadmap/}\allowbreak\file{protocol-blueprint.md:1130}) is therefore the load-bearing companion of the
  master theorem, and it does not exist yet.
\item \textbf{Positive control.} \lean{checkBit} (\lean{verify := fun s _}\allowbreak\lean{ => if s then pure () else failure})
  is knowledge sound at error $0$ with the trivial extractor (\lean{checkBit_sound_at_zero}): the
  intended object inhabits the definition.
\item \textbf{Constant state functions.} For any knowledge state function and any statement with
  no witness, the state at round $0$ is false for every witness (\lean{stateFunction_false_at});
  for a statement with a valid witness it is true (\lean{stateFunction_true_at}). So a knowledge
  state function is never constantly true nor constantly false as soon as the relation has a
  statement with a witness and one without.
\item \textbf{Trivial output relation.} \lean{full_iff_of_univ}: for a guarded verifier and
  \lean{relOut = Set.univ}, the last law of a knowledge state function is equivalent to \enquote{if
  the check passes on $(s, \mathit{tr})$ then the state at the full transcript holds at the
  extracted witness}. With the spine's \lean{Seam.done I} (\lean{fun _ _ ↦ True},
  \file{LeanerVM/}\allowbreak\file{Protocol/}\allowbreak\file{Spine/}\allowbreak\file{Seams.lean:191}) and its output statement \lean{Unit},
  \enquote{accept} means \enquote{the verifier returns rather than fails}, and law 3 says exactly
  \enquote{accepted $\Rightarrow$ state true}. A design consequence: since the final output
  statement is \lean{Unit}, the last phase can only reject by failing; a verifier that returned a
  Boolean verdict would count \lean{false} as an acceptance under \lean{Set.univ}.
\end{itemize}

\subsection{The oracle verifier is judged through \texorpdfstring{\texttt{toVerifier}}{toVerifier}}\label{sec:gen-lib-ark-toverifier}

Every security predicate ArkLib gives an \lean{OracleVerifier} (\lean{OracleVerifier.}\allowbreak\lean{rbrSoundness},
\lean{rbrKnowledgeSoundness}, \lean{rbrKnowledgeSoundnessWith},
\file{Security/}\allowbreak\file{RoundByRound.lean:729-758}) is the predicate of its plain image \lean{toVerifier}
(\cref{sec:gen-lib-ark-reductions}), and leanerVM states its master theorem on
\lean{(leanVmVerifier P).}\allowbreak\lean{toVerifier}. \lean{simOracle2} answers each query of the oracle verifier
from the actual oracle statements and the actual messages of the transcript, and the output
oracles are materialised. So the plain verifier's verdict on a transcript is the oracle verifier's
verdict when its queries are answered faithfully; the state function and the extractor are
functions of the whole transcript, including the full contents of every oracle message (the stack
$q$, $2^{\mu}$ field elements). This is the standard reading of IOP soundness, and nothing is lost
for soundness at the oracle-protocol level. Three things are outside the statement:
\begin{itemize}
\item the verifier's query complexity and locality (an \lean{OracleVerifier} may query every
  point; \lean{OracleVerifier.numQueries} is a \lean{sorry} definition, \file{Basic.lean:505-507});
  they matter for the cost of the compiled protocol, not for its soundness;
\item the \emph{type} of the oracle message: a prover in the security game can only send an element
  of \lean{Column I.μ} (a table of $2^{\mu}$ values in $\K$), that is, a genuine multilinear over
  the base field of the announced size. Whether a prover of the compiled protocol can commit to
  something else is the commitment scheme's problem (extractability or list binding of WHIR,
  Layer 11), not this theorem's;
\item the composition \lean{OracleVerifier.append} (\file{Append/}\allowbreak\file{Basic.lean:619-629}) routes the
  second verifier's queries through the first's output-oracle simulation;
  \lean{append_toVerifier} (\file{:658}, proved) says its \lean{toVerifier} is the plain
  \lean{Verifier.append} of the two \lean{toVerifier}s, which is what lets leanerVM compose state
  functions of plain verifiers.
\end{itemize}

\subsection{Must an extractor compute?}\label{sec:gen-lib-ark-extractors}

\paragraph{A classical-choice extractor is accepted by the type.} ArkLib forces nothing:
\lean{Extractor.RoundByRound} is a structure of Lean functions with no computability requirement,
and ArkLib itself defines a \lean{noncomputable} one with \lean{h.choose}
(\lean{Extractor.}\allowbreak\lean{RoundByRoundOneShot.}\allowbreak\lean{toRoundByRoundOfRel}, \file{Security/}\allowbreak\file{RoundByRound.lean:118-123}).
The probe \code{Extractors} shows the consequence: \lean{chooser R} (\lean{extractOut} chooses a
witness of \lean{(s, ·) ∈ R} if one exists, classically) is accepted as an extractor, and
\lean{chooser_sound} proves \lean{rbrKnowledgeSoundnessWorstCaseWith … (chooser R) … (fun _}\allowbreak\lean{ => 0)}
for the verifier that decides the language of \lean{R}, for \emph{every} relation \lean{R} with
that language. Knowledge has been proved of nothing.

\paragraph{The existential form is soundness only.} In the existential forms
(\lean{rbrKnowledgeSoundness}, \lean{rbrKnowledgeSoundnessWorstCase}, leanerVM's
\lean{piop_}\allowbreak\lean{rbrKnowledgeSoundness_}\allowbreak\lean{exists}) the extractor is chosen by the prover of the theorem, so
such a theorem proves soundness, not knowledge (\cref{f:lib-existential-rbr}).

\paragraph{The named form is meaningful with the extractor read.} What tells a computable
extractor:
\begin{itemize}
\item \texttt{\#print axioms} does not tell: \lean{piopExtractor}, \lean{commitExtractor} and
  \lean{PublicInput.extractor} all report \lean{Classical.choice} (probe \code{AxiomsLeanerVM}),
  through the proofs and instances inside their types, as does nearly every declaration built on
  Mathlib.
\item What tells is the compiler: a \lean{def} without the \lean{noncomputable} keyword is
  accepted only if Lean can generate code for it, and Lean refuses a definition using classical
  choice in a value position (\enquote{failed to compile definition, consider marking it as
  'noncomputable' because it depends on 'Classical.propDecidable'}, probe
  \code{ExtractorsExpectedFailure}, line 12). The refusal is transitive: a compiled definition
  cannot call a \lean{noncomputable} one.
\item \texttt{\#eval} is the positive test: it runs the code.
\end{itemize}
Applied to leanerVM (probe \code{Extractors}): \lean{commitExtractor I},
\lean{PublicInput.extractor I}, \lean{piopExtractor P S}, \lean{(commitSecurity I).}\allowbreak\lean{extractor},
\lean{(publicInputSecurity I).}\allowbreak\lean{extractor} and an appended extractor are each the body of a compiled
\lean{def}; the commit extractor evaluated on the toy instance returns the honest stack
(\lean{true}, twice), and the public-input extractor returns \lean{()}. The definitions are
moreover one line each: \enquote{read the message} (\file{SendOracle.lean:137-138}) and
\enquote{return ()} (\file{PublicInput.lean:320-321}), so their efficiency is read off directly,
which matters because Lean-computable is not polynomial-time: for a finite witness type and a
decidable relation, exhaustive search is a computable extractor for any sound verifier.

Two limits of the check:
\begin{enumerate}
\item \lean{piopExtractor P S} is computable \emph{as a function of} the bundle
  \lean{S : P.Security}, whose \lean{extractor} field any future phase fills. The check must be
  repeated on each phase's \lean{Security} value (\lean{publicInputSecurity} passes today). The
  type does not enforce it.
\item \lean{Component.Def} carries the real-valued error \lean{err : … → ℝ≥0}, so every phase
  definition with a non-zero error is \lean{noncomputable} (\lean{publicInputPhase},
  \file{LeanerVM/}\allowbreak\file{Protocol/}\allowbreak\file{PublicInput.lean:374-381}, \enquote{noncomputable because of the error
  alone}). Lean then refuses to compile anything that takes the phase's definition as an argument:
  the probe \code{ExtractorsExpectedFailure} shows \lean{(publicInputPhase I).}\allowbreak\lean{red.}\allowbreak\lean{verifier} and
  \lean{extractorOf (publicInputPhase I) (publicInputSecurity I)} rejected, although the verifier
  and the extractor are compiled when named directly. \lean{leanVmVerifier P := P.}\allowbreak\lean{toDef.}\allowbreak\lean{red.}\allowbreak\lean{verifier}
  and \lean{piopExtractor P S} take \lean{P} explicitly, so on the real bundle of phases neither
  will be a compiled definition. This does not affect the theorems; it affects the blueprint's
  promises that the extractor and the honest prover \enquote{compute} and the planned
  code-generation probes (\cref{f:lib-err-noncomputable}).
\end{enumerate}

\subsection{Perfect completeness, and completeness with an error}\label{sec:gen-lib-ark-completeness}

\begin{define}{Completeness (\lean{Reduction.completeness}, \lean{perfectCompleteness})}
For every $(\mathit{stmtIn}, \mathit{witIn})$ in the input relation, run the honest prover and the
verifier (\lean{Reduction.run}, \file{Execution.lean:209-216}: the prover runs, then the verifier
on the transcript; challenges answered uniformly by \lean{challengeQueryImpl}, shared oracle by
\lean{impl} from a state drawn by \lean{init}); the probability that the run does not fail, the
verifier's output statement with the prover's output witness is in the output relation,
\textbf{and} the prover's output statement equals the verifier's, is at least $1 - \varepsilon$.
Perfect completeness is $\varepsilon = 0$, so probability $1$
(\lean{perfectCompleteness_}\allowbreak\lean{eq_prob_one}, \file{:165}). Failure (the verifier rejecting) counts
against completeness. \lean{OracleReduction.}\allowbreak\lean{perfectCompleteness} (\file{:469}) is the same for
\lean{toReduction}.
\end{define}
\begin{leancode}
def completeness (relIn : Set (StmtIn × WitIn))
    (relOut : Set (StmtOut × WitOut))
    (reduction : Reduction oSpec StmtIn WitIn StmtOut WitOut pSpec)
    (completenessError : ℝ≥0) : Prop :=
  ∀ stmtIn : StmtIn,
  ∀ witIn : WitIn,
  (stmtIn, witIn) ∈ relIn →
    let pImpl : QueryImpl (oSpec + [pSpec.Challenge]ₒ) (StateT σ ProbComp) :=
      QueryImpl.addLift impl challengeQueryImpl
    Pr[fun ⟨⟨_, (prvStmtOut, witOut)⟩, stmtOut⟩ =>
        ((stmtOut, witOut) ∈ relOut ∧ prvStmtOut = stmtOut) | OptionT.mk do
          (simulateQ pImpl (reduction.run stmtIn witIn).run).run' (← init)] ≥ 1 - completenessError

/-- A reduction satisfies **perfect completeness** if it satisfies completeness with error `0`. -/
def perfectCompleteness (relIn : Set (StmtIn × WitIn)) (relOut : Set (StmtOut × WitOut))
    (reduction : Reduction oSpec StmtIn WitIn StmtOut WitOut pSpec) : Prop :=
  completeness init impl relIn relOut reduction 0
\end{leancode}
\src{\file{ArkLib/}\allowbreak\file{OracleReduction/}\allowbreak\file{Security/}\allowbreak\file{Basic.lean:89-105}, ArkLib \code{dca90385}}

\paragraph{What the composition theorem needs.} leanerVM's \lean{Component.Complete.append}
(\file{Component.lean:162-169}) applies the following theorem of ArkLib,
\lean{OracleReduction.}\allowbreak\lean{append_}\allowbreak\lean{perfectCompleteness_}\allowbreak\lean{of_guarded_}\allowbreak\lean{verifiers}, which is proved (probe
\code{AxiomsArkLibBuilt}):
\begin{leancode}
/-- Oracle reductions with guarded ordinary verifiers compose perfectly when prover execution
factors at the seam and the suffix is perfectly complete from every shared state. -/
theorem append_perfectCompleteness_of_guarded_verifiers
    (R₁ : OracleReduction oSpec Stmt₁ OStmt₁ Wit₁ Stmt₂ OStmt₂ Wit₂ pSpec₁)
    (R₂ : OracleReduction oSpec Stmt₂ OStmt₂ Wit₂ Stmt₃ OStmt₃ Wit₃ pSpec₂)
    (V₁ : R₁.toReduction.verifier.GuardedForm) (V₂ : R₂.toReduction.verifier.GuardedForm)
    (hSeam : ∀ hn : 0 < n,
      R₁.prover.OutputIsPure ∨ pSpec₂.dir ⟨0, hn⟩ = .P_to_V)
    (h₁ : R₁.perfectCompleteness init impl rel₁ rel₂)
    (h₂ : ∀ s : σ, R₂.perfectCompleteness (pure s) impl rel₂ rel₃) :
    (R₁.append R₂).perfectCompleteness init impl rel₁ rel₃ := by
  change (R₁.append R₂).completeness init impl rel₁ rel₃ 0
  simpa only [add_zero] using
    append_completeness_of_guarded_verifiers R₁ R₂ V₁ V₂ hSeam h₁ h₂
\end{leancode}
\src{\file{ArkLib/}\allowbreak\file{OracleReduction/}\allowbreak\file{Composition/}\allowbreak\file{Sequential/}\allowbreak\file{OracleCompleteness.lean:54-67}, ArkLib \code{dca90385}}

Its hypotheses: guarded forms of both plain verifiers (so that a run is \enquote{the prover's run,
then a deterministic check and verdict}, \file{GuardedCompleteness.lean:44-48}); the seam condition
\lean{hSeam} (if the second protocol has a round: the first prover's output is a pure function of
its state, \lean{Prover.OutputIsPure}, \emph{or} the second protocol starts with a prover message),
which is what makes the appended prover's execution factor into the two provers' executions
(\lean{Prover.}\allowbreak\lean{SimulatedAppendFactorization}, \file{Append/}\allowbreak\file{Simulation.lean:131-153}); completeness
of the first component from \lean{init}; and completeness of the second \textbf{from every
deterministic state \lean{pure s}} of the shared oracle, because the second component starts in
whatever state the first left. leanerVM asks every component for completeness from every
\lean{init} (\lean{Complete.complete}, \file{Component.lean:84-85}), which gives both. Over the
empty shared oracle the guarded forms and the pure outputs are automatic
(\lean{Verifier.}\allowbreak\lean{GuardedForm.}\allowbreak\lean{ofEmpty}, \lean{Prover.}\allowbreak\lean{instOutputIsPureEmpty};
\cref{sec:gen-lib-ark-port}).

\paragraph{Completeness with an error.} ArkLib at the old pin composes completeness \textbf{with
errors}, and the errors add: \lean{Reduction.}\allowbreak\lean{append_}\allowbreak\lean{completeness_}\allowbreak\lean{of_guarded_}\allowbreak\lean{verifiers}
(\file{GuardedCompleteness.lean:161-173}, error $\varepsilon_1 + \varepsilon_2$),
\lean{OracleReduction.}\allowbreak\lean{append_}\allowbreak\lean{completeness_}\allowbreak\lean{of_guarded_}\allowbreak\lean{verifiers} (\file{OracleCompleteness.lean:39-52}),
and the $n$-ary \lean{seqCompose_}\allowbreak\lean{completeness_}\allowbreak\lean{of_guarded_}\allowbreak\lean{verifiers} (\file{GuardedNary.lean:45-59},
\file{OracleCompleteness.lean:77-99}, error $\sum_i \varepsilon_i$); all proved (probe
\code{AxiomsArkLibBuilt}). The perfect versions are their $\varepsilon = 0$ corollaries. So a
phase complete only up to an error can be composed in ArkLib today.

leanerVM's spine cannot take one: \lean{Component.}\allowbreak\lean{Complete.}\allowbreak\lean{complete} is \lean{perfectCompleteness}
(\file{Component.lean:84-85}), \lean{Complete.append} uses the perfect theorem, and the
blueprint's \lean{FlockInterface} (\file{docs/}\allowbreak\file{roadmap/}\allowbreak\file{protocol-blueprint.md:1083}) demands
\lean{perfectCompleteness : reduction.}\allowbreak\lean{perfectCompleteness …}. The blueprint's list of wrong
readings, item 20 (\file{:1323-1325}), says: \enquote{Perfect completeness needs no
exceptional-challenge clause in Layers 4--8 and 10: no inverse of a challenge is taken. Flock's
$(1 + r_{\mathrm{eq}})^{-1}$ at $r_{\mathrm{eq}} = 1$ is \#3's, inside \lean{flockError}. Witness:
\lean{piop_perfectCompleteness} has no side condition on challenges.} The pinned Rust Flock prover
does take that inverse (\file{crates/}\allowbreak\file{flock/}\allowbreak\file{src/}\allowbreak\file{zerocheck.rs:117},
\code{let g0 = (claim + r_eq * g1) * (F192::ONE + r_eq).inv();}), which is undefined at
$r_{\mathrm{eq}} = 1$ in characteristic 2 (the Rust \code{inv} returns $0$ there,
\file{crates/}\allowbreak\file{primitives/}\allowbreak\file{src/}\allowbreak\file{field/}\allowbreak\file{gf2_64x3.rs:144-153}, so the honest prover sends a wrong
$g_0$). Whether the honest Flock prover then fails is for the Flock review; if it does, the Flock
phase has completeness error $1/\abs{\E}$ per such challenge, not $0$, and a completeness failure
cannot be charged to \lean{flockError}, which is a soundness error. Then \lean{FlockInterface}'s
\lean{perfectCompleteness} field is unsatisfiable, \lean{Phases.Complete} cannot be built for the
real protocol, and \lean{piop_perfectCompleteness} is vacuous for it. The fix is available at the
pin: give \lean{Component.Complete} an error field and use the additive theorems
(\cref{f:lib-completeness-error}).

\par\endgroup